\section{Introduction}
\label{sec:intro}

A robot maps a room and returns after a chair has moved. The earlier depth images
were correct when they were captured, but some now describe an obsolete scene.
Keeping all of them leaves the chair in two places. Starting over loses parts
of the room that the robot does not revisit. A useful update must answer two
questions: which observations still describe the current scene, and where should
the previous map remain in use?

\begin{figure*}[t]
\centering
\begingroup
\setlength{\fboxsep}{6pt}
\newcommand{\updatebox}[2]{\fbox{\begin{minipage}[c][2.1cm][c]{0.26\textwidth}
\centering\small\textbf{#1}\par\vspace{5pt}#2\end{minipage}}}
\updatebox{Previous map + new observations}{Stored geometry and object history\\
Current RGB-D, poses, and identities}
\hspace{2pt}$\longrightarrow$\hspace{2pt}
\updatebox{Select temporally valid evidence}{Which observations still apply?\\
Where should memory remain?}
\hspace{2pt}$\longrightarrow$\hspace{2pt}
\updatebox{Update the native map}{TSDF: fuse the selected evidence\\
Extract current surfaces; retain history}
\endgroup
\vspace{4pt}
\caption{\textbf{Select evidence before integrating geometry.} A state change
ends the use of observations from the old state. Valid current measurements
replace prior influence where they provide evidence; eligible memory supplies
unobserved regions. Temporal selection and geometric integration have distinct
roles, with TSDF fusion providing the concrete realization developed here.}
\label{fig:update_overview}
\end{figure*}


Dynamic mapping already addresses motion and changes during a recording.
Object-level systems reconstruct moving entities with separate surface or volume
models~\cite{runz2017cofusion,xu2019midfusion,strecke2019emfusion}, while joint
estimators recover camera and object motion~\cite{bescos2021dynaslam2,morris2025dynosam}.
Changing-SLAM maintains object persistence across observation gaps
~\cite{soares2023changing}; Khronos combines visible motion with temporal
reconciliation~\cite{schmid2024khronos}; and GaME updates Gaussian geometry
after out-of-view changes~\cite{yugay2026game}. SuperMap maintains online
object maps with open-vocabulary instances and temporal relations
~\cite{zhao2026supermap}.

Long-term mapping carries information between visits and revises an earlier map
~\cite{pomerleau2014longterm,lazaro2018efficient}. ELite and ProbPer-LiLo model
which content persists over repeated traversals~\cite{gil2025elite,ali2026probperlilo}.
Panoptic Multi-TSDFs distinguishes persistent, absent, and unobserved submaps
~\cite{schmid2022panoptic}. POCD and POV-SLAM reason about object changes
~\cite{qian2022pocd,qian2023povslam}, ObVi-SLAM transfers object landmarks
~\cite{adkins2024obvislam}, and OASIS-Map associates objects across sessions
~\cite{oh2026oasismap}. These approaches motivate a shared question beneath
particular map representations: when should an earlier measurement continue to
influence the current map?

We address this question by separating \emph{temporal evidence selection} from
\emph{geometric integration}. Our first rule is that an observation contributes
to current geometry only while the scene state it measured remains valid. If an
object moves, measurements of its previous state belong to history rather than
to the current reconstruction. Our second rule selects between current evidence
and memory. Where the current run has valid depth evidence, that evidence defines
the update. Where it has none, the previous map is retained, subject to the
current object-state assignment. A ray that passes through an old surface is
current evidence even when it produces no replacement surface there.

The distinction matters before geometry is converted into a mesh. A mesh records
surfaces, whereas the depth observations used to build it also identify free
space. Comparing two final meshes alone loses this distinction between a region
that was observed empty and one that was not observed. We instead use temporal
state and ray evidence to select the observations that enter reconstruction.
For a TSDF map, the resulting field is a weighted average of currently valid
observations and a spatially gated prior. Object and background surfaces are
extracted after this selection, using the same update rule.

The temporal policy is defined independently of a particular mapper. A mapper
supplies geometry queries and native integration operations; the policy decides
which observations and inherited states remain eligible. We develop this
separation through a TSDF formulation and a Khronos-based implementation. This
preserves the backbone's motion and object-state estimates while moving the
persistence decision into the evidence used to construct the map.

Our experiments examine three consequences of persistent updates: whether old
states are removed, whether previous coverage remains useful, and whether the
new run continues to handle its own changes. Repeated Kinect and Isaac Sim
recordings exercise the complete pipeline. 3RScan supplies public revisit cases,
while Bonn and TUM test single-run object motion and moving people.

Our contributions are:
\begin{itemize}
    \item a temporal update formulation that selects observations according to
    their state validity and selects prior geometry according to current
    measurement coverage;
    \item a TSDF realization that applies these selections before object and
    background surface extraction, with object identities and historical states
    maintained by the same temporal reasoning;
    \item a repeated-mapping evaluation covering changed-state accuracy,
    obsolete-surface removal, geometry coverage, and single-run behavior.
\end{itemize}
